% !TEX root = cubic_artin_windows.tex
\documentclass[a4paper,11pt]{ltjsarticle}
\usepackage[margin=25mm,headheight=16pt]{geometry}
\usepackage{amsmath,amssymb,amsthm,mathtools}
\usepackage{enumitem,booktabs}
\usepackage[unicode,hidelinks]{hyperref}
\usepackage{fancyhdr}
\usepackage{microtype}
\setlength{\emergencystretch}{2em}
\allowdisplaybreaks[2]
\setlist{itemsep=0.2em,topsep=0.4em}
\pagestyle{fancy}\fancyhf{}
\fancyhead[L]{\small 有限Artin基底上の残余圏と行列因子分解}
\fancyhead[R]{\small 2026-11-13 / v1}
\fancyfoot[C]{\thepage}
\newcommand{\CC}{\mathbb C}
\newcommand{\ZZ}{\mathbb Z}
\newcommand{\PP}{\mathbb P}
\newcommand{\MF}{\operatorname{MF}}
\newcommand{\Ku}{\operatorname{Ku}}
\newcommand{\Id}{\operatorname{id}}
\newtheorem{theorem}{定理}[section]
\newtheorem{proposition}[theorem]{命題}
\newtheorem{lemma}[theorem]{補題}
\newtheorem{corollary}[theorem]{系}
\theoremstyle{definition}
\newtheorem{definition}[theorem]{定義}
\newtheorem{example}[theorem]{例}
\theoremstyle{remark}
\newtheorem{remark}[theorem]{注意}
\renewcommand{\proofname}{証明}
\numberwithin{equation}{section}

\hypersetup{pdftitle={有限Artin基底上の残余圏と行列因子分解}}
\title{有限Artin基底上の残余圏と行列因子分解\\
\large 相対window、底変更、完備化における生成性の障害}
\author{}
\date{2026年11月13日／研究草稿 v1}
\begin{document}
\maketitle
\begin{abstract}
$\mathbb C[h]/(h^{N+1})$上の次数別有限自由な中心三次対で、
特殊ファイバーが多項式環と滑らかな三次式となるものを考える。
任意の$N\ge0$と変数数$d\ge5$について、線束生成残余圏を
graded matrix factorizationと比較し、底変更・rotation・指定三次周期との両立を示す。
正則局所基底からのArtin底変更も記述する。
他方、全Artin切断が相対smoothであっても、完備基底でのsmoothnessは従わないことを、
射影直線上の生成対象の変形と四次元strict dg代数により示す。
MF同値から可換残余圏へ移す際には、MF同値そのものの存在を別の入力として区別する。
\end{abstract}

\section*{問題設定と証明の構成}
本稿は任意の局所Artin環への一般化ではなく、一変数の全有限切断
$R_N=\mathbb C[h]/(h^{N+1})$を扱う。
$d$は変数数、超曲面の次元は$d-2\ge3$である。
線束で生成される部分を扱い、非正則な基底上で全coherent圏を
通常の絶対特異圏と比較することはしない。

中心結果は定理\ref{am:thm:comparison}の
$\mathcal K_B^{\rm lin}\simeq\MF^{\ZZ}(A,\Omega)$である。
余核による相対stabilization、線束とMFのHomの有限性・底変更、
Ind圏での生成子の持上げを順に用いる。
この構成の加群論的入力はChen\cite{AMChen}、射影化の入力は
Artin--Zhang\cite{LP-AZ}とBurke--Stevenson\cite{LP-BS}である。
既知の三角圏同値をそのままdg同値と呼ばず、exact圏とdg商による比較を記す。
この組合せの先行研究上の初出・包括的新規性は未確認である。

第\ref{ar:sec:regular}節は正則局所模型からの底変更に必要な証明を収める。
第\ref{ar:sec:application}節では、可換MFとの同値を別途仮定して残余圏へ移す。
可換三次式やそのMF同値を構成する定理ではない。
第\ref{as:sec:obstruction}節は、有限Artin基底での生成性と
完備基底での生成性を同一視できない理由を示す。
その出発点となる生成対象の変形例は既知である\cite{AMSasha}。

% Current finite-Artin comparison; extracted with complete proofs.
\section{Artin基底における線束生成残余圏と行列因子分解}
\label{am:sec:comparison}

本節では絶対特異圏を用いず、周囲の代数上perfectな複体だけを用いて、
有限階数graded MFと線束から生成される射影残余圏を比較する。
ここで構成するのは射影残余圏とMF圏の比較であり、
非可換MF圏と可換MF圏の間の別の同値の存在は仮定しない。
非正則Artin基底上で通常の有界coherent圏全体へOrlovの定理を適用することはしない。
相対的なstabilizationの加群論的入力は
Chen\cite[Theorem 2.10, Corollary 2.11, Theorem 7.2]{AMChen}である。
同論文の定理は三角圏についてのものであるため、以下ではexact圏とdg商を通じて
実際のenhanced関手を構成する。
本節の組合せと同じ一般性での先行研究上の初出は未確認である。

\subsection{入力と相対的な特異圏}

変数の数を$d\ge5$、Artin切断の次数を$N\ge0$とし、
\[
 R_N=\CC[h]/(h^{N+1}),\qquad k=R_N/(h),\qquad
 a=d-3,\qquad i=1-a=4-d
\]
と置く。$A=\bigoplus_{j\ge0}A_j$は中央$R_N$-代数で、$A_0=R_N$、
各$A_j$は有限自由$R_N$-加群、
$\overline A=A\otimes_{R_N}k\simeq k[z_1,\ldots,z_d]$と仮定する。
$\Omega\in Z(A)_3$の剰余$F$は滑らかな射影cubicを定め、$B=A/(\Omega)$とする。
以下の構成は内部次数を保つ。本節のMFと厚いdg部分圏は冪等完備化して扱う。

この入力から、左右noetherian性、$R_N$の左右$A$-加群としての有限自由Koszul型分解、
$\Omega$の正則性、$B_j$の自由性が従う。
これらは正則基底上の有限global dimensionと区別する。
実際、$h$による有限filtrationの随伴次数環がnoetherianなので$A,B$はnoetherianである。
この持上げは、基底の正則性ではなく次数別自由性を使う。
剰余の自由分解の零次項の生成元の像を同じ次数に持ち上げると、
次数別Nakayamaにより全射を得る。
各次数で有限自由加群への全射なので分裂し、核も次数別有限自由である。
この核の剰余を正しく計算できるため、次の自由項を同様に持ち上げて帰納する。
最後の核は剰余が零であり、再び次数別Nakayamaにより零となる。
従って多項式環のKoszul分解から
\begin{equation}\label{am:eq:relative-koszul}
 P_j=A(-j)^{\binom dj}\quad(0\le j\le d)
\end{equation}
という分解を得る。各内部次数で$R_N$-分裂し、全てのArtin商への底変更後も完全である。
$\Omega:A_{j-3}\to A_j$は剰余上で単射なので、有限自由加群の行列の単元小行列式を用いて
$R_N$上でも分裂単射となる。従って$B_j$は自由である。
この議論から$\operatorname{gldim}A<\infty$は結論しない。

\begin{definition}[周囲でperfectな特異圏]\label{am:def:relative-singularity}
標準dg enhancement内で
\[
 \mathcal D_{B/A}
 =\{M\in D^b(\operatorname{gr}B):
       \operatorname{Res}_A M\in\operatorname{Perf}(\operatorname{gr}A)\}
\]
を取り、
\[
 \mathcal S_{B/A}
 =\bigl(\mathcal D_{B/A}/\operatorname{Perf}(\operatorname{gr}B)\bigr)^{\mathrm{idem}}
\]
をそのdg商の冪等完備化とする。
$\operatorname{Perf}(\operatorname{gr}B)\subset\mathcal D_{B/A}$は
$0\to A(-3)\xrightarrow{\Omega}A\to B\to0$から従う。
\end{definition}

ここでのperfect性は周囲の\emph{$A$上}のperfect性である。
$B(j)$やその内部切断は、$R_N$上では一般に無限階数であり、
$\operatorname{Perf}(R_N)$の対象だとは述べない。
また$\mathcal S_{B/A}$を絶対特異圏$D_{\rm sg}(B)$と同一視しない。

\begin{lemma}[相対dg stabilization]\label{am:lem:stabilization}
上の入力に対し、内部shiftと三次周期性を保つdg Morita同値
\begin{equation}\label{am:eq:stabilization}
 \operatorname{cok}:
 \MF^{\ZZ}(A,\Omega)\xrightarrow{\sim}\mathcal S_{B/A}
\end{equation}
が存在する。そのMorita逆を$\operatorname{stab}$と書く。
この同値は有限自由MFの余核とその周期自由分解から誘導される。
Artin商への底変更に関して、これらの構成は自然である。
\end{lemma}

\begin{proof}
まず、次数付き有限射影$A$-加群は次数付き有限自由である。
実際、そのaugmentation商は局所環$R_N$上の次数付き有限射影加群なので自由であり、
その基底を持ち上げた自由加群からの写像に次数付きNakayamaを適用すればよい。

$\mathcal E$を、有限生成次数付きGorenstein-projective $B$-加群のうち
$A$上有限射影次元を持つもののexact圏とする。
Chenのchange-of-rings定理により、非零のこのような加群は
$A$上の射影次元が一である。
正則中心元の場合の同論文の証明は、全ての射と分解を斉次に取れば次数付き圏でも成立する。
特に$\mathcal E$はFrobenius exact圏であり、そのprojective-injective対象は
有限次数付き射影$B$-加群である。
syzygyとcosyzygyで有限$A$-射影次元が保たれること、およびその次元が再び一になることが、
このFrobenius部分圏への制限を保証する。

$\mathcal F$を有限自由graded MFのexact圏とする。
exact列は二つの成分で$A$-分裂する列であり、projective-injective対象は
$(1,\Omega)$、$(\Omega,1)$型の可縮factorizationの直和因子である。
余核はexact関手
\[
 \operatorname{cok}:\mathcal F\longrightarrow\mathcal E
\]
を与える。余核加群の$A$-自由分解の長さは一である。
逆に$M\in\mathcal E$の斉次分解
$0\to P_1\xrightarrow{\varphi}P_0\to M\to0$を取ると、
$\Omega M=0$より$\Omega\Id_{P_0}$は$\varphi$を通って一意にfactorする。
$\varphi$の単射性から他方の積も$\Omega\Id_{P_1}$となり、MFを得る。
射の持上げとhomotopyの同定を含め、これはChenのTheorem 7.2の余核同値である。
従って余核関手はstable categoryの三角同値を与える。

この三角同値だけでdg同値を宣言せず、次のdg商を取る。
exact関手から実際のdg関手
\begin{equation}\label{am:eq:exact-dg-quotients}
 \frac{D^b_{\rm dg}(\mathcal F)}{K^b_{\rm dg}(\operatorname{proj\!-\!inj}\mathcal F)}
 \longrightarrow
 \frac{D^b_{\rm dg}(\mathcal E)}{K^b_{\rm dg}(\operatorname{proj}B)}
\end{equation}
を得る。Frobenius exact圏の右辺・左辺の$H^0$は対応するstable categoryである。
これはprojective-injective分解を用い、有限複体をその高いsyzygyへ移すことで確認できる。
従って\eqref{am:eq:exact-dg-quotients}は$H^0$で同値であり、
shiftを含む任意の対象対へ同じ事実を適用すると全Hom cohomologyで同型となる。

左側のdg商から元のMF dg圏へは、有限複体の有限totalizationを取る。
成分別split exactな複体のtotalizationは可縮である。
具体的に、横微分$d$と横contraction$s$を
$ds+sd=1$、$s^2=0$とし、符号を含む縦MF微分を$\delta$とする。
$d\delta+\delta d=0$、$\delta^2=w$で、$w$は中央曲率である。
$s$は横次数を下げるので、横方向の有限性から
$1+\delta s$、$1+s\delta$は有限和で逆を持つ。
\[
 q=s(1+\delta s)^{-1}=(1+s\delta)^{-1}s
\]
と置けば、直接展開により
\[
 (1+s\delta)\bigl(1-(d+\delta)q-q(d+\delta)\bigr)
 (1+\delta s)=-sws=0
\]
となる。最後は中央性と$s^2=0$を用いた。
従って$(d+\delta)q+q(d+\delta)=1$であり、
totalizationは実際に可縮である。この計算は曲率の冪零性を必要としない。
可縮MFも零へ送られるため、totalizationは左側のdg商を通る。
$H^0$では通常のMF homotopy categoryへの同値であり、同様に全Homで準同型となる。

右側から$\mathcal S_{B/A}$へのdg関手は、$\mathcal E$の$B$-加群圏への包含から得る。
完全射影分解を用いると、Gorenstein-projective加群間のstable Homは特異圏のHomに一致する。
同じ分解のshiftを用いることで全Hom cohomologyについて一致する。
本質的全射性も具体的に確認できる。
$M\in\mathcal D_{B/A}$の有界上の有限生成$B$-射影分解を十分低い位置で切断すると、
特異圏で$M$はある$B$-加群$L$のshiftに同型となり、$L$は$A$上有限射影次元を持つ。
ChenのTheorem 2.10より$L$は$B$上有限Gorenstein-projective次元を持つので、
さらに高い$B$-syzygyを取れば$\mathcal E$に属する。
これで右側のdg商と$\mathcal S_{B/A}$の同値を得る。
三つの実際のdg関手のzigzagにより\eqref{am:eq:stabilization}が得られる。

最後に底変更を確認する。
MFの余核$M$の各内部次数の分解は有限自由$R_N$-加群間の列である。
剰余上で最初の射は単射である。実際、二つのfactorization mapの積は
非零多項式$F$の乗法なので、どちらのmapも単射である。
単元小行列式により各次数の余核$M_j$は$R_N$-自由となる。
従って余核、exact列、有限totalizationはArtin商への底変更と両立する。
周期分解とdg商から得た比較は、Morita localizationの普遍性を通じて自然となる。
選択したcontractionそのものが全ての底変更と厳密に可換することは必要としない。
内部shiftと周期性は全ての段階で保たれ、Morita逆もこれらと両立するように取れる。
\end{proof}

\subsection{Artin基底上の線束Homと有限window}

$\pi:\operatorname{Gr}B\to\operatorname{QGr}B$を射影商とし、
$\mathcal O(j)=\pi B(j)$と置く。
導来sectionを$\mathbf R\omega$と書く。
この段階では$D(\operatorname{QGr}B)$内で線束から生成するdg部分圏を用い、
全ての有界coherent複体についての底変更定理は要求しない。

\begin{lemma}[線束のsectionと底変更]\label{am:lem:line-bundles}
自然なderived Homの底変更は各Artin商と両立し、
\begin{equation}\label{am:eq:line-hom}
 \operatorname{Ext}^q(\mathcal O(r),\mathcal O(s))=
 \begin{cases}
 B_{s-r}&q=0,\\
 B_{r-s-a}^{\vee}&q=d-2,\\
 0&\text{その他}
 \end{cases}
\end{equation}
である。ここで$(-)^\vee=\operatorname{Hom}_{R_N}(-,R_N)$、負添字は零とする。
また$0\le j\le d-1$について自然に
\begin{equation}\label{am:eq:window-sections}
 (\mathbf R\omega\mathcal O(j))_{\ge i}\simeq B(j)_{\ge i}
\end{equation}
となり、この同定もArtin商への底変更と両立する。
\end{lemma}

\begin{proof}
式\eqref{am:eq:relative-koszul}を$A$で双対化する。
各内部次数は有界有限自由$R_N$-複体であり、剰余は通常のKoszul双対性を満たす。
単元成分の消去によって
\[
 \mathbf R\underline{\operatorname{Hom}}_A(R_N,A)
 \simeq R_N(d)[-d]
\]
となる。中心正則元のchange-of-ringsから
\[
 \mathbf R\underline{\operatorname{Hom}}_B(R_N,B)
 \simeq R_N(a)[-(d-1)]
\]
を得る。この二つの計算は基底の有限global dimensionを用いない。

標準の導来section--torsion三角を用いる。
ここで必要な$\chi$はArtin基底でも確認できる。
$N>0$では$B/(h)$の$B$-自由分解は、中央の$h,h^N$による周期分解である。
有限生成$B$-加群$M$について、そのHom cohomologyは有限生成$B/(h)$-加群となる。
change-of-rings spectral sequenceと可換noetherian環$B/(h)$の$\chi$から
$\operatorname{Ext}_B^q(k,M)$は各$q$で内部次数について有限となる。
$R_N$の有限な$k$-filtrationにより、$\operatorname{Ext}_B^q(R_N,M)$も同じ性質を持つ。
$N=0$では通常の可換計算である。
従ってArtin--Zhangのsectionとlocal cohomologyの枠組を適用できる
\cite{LP-AZ}。

$M_\ell=B/B_{\ge\ell}$の内部次数filtrationの商は
$B_j\otimes_{R_N}R_N(-j)$である。
直前のrelative Gorenstein計算と$B_j$の自由性から
\[
 H^q(\mathbf R\tau B)_t=
 \begin{cases}
 B_{-t-a}^{\vee}&q=d-1,\\
 0&q\ne d-1
 \end{cases}
\]
を得る。固定した$t$では$j=-t-a$の一層のみが寄与し、$\ell>j$で安定する。
この計算では、次数filtrationの各商の導来Homがcohomological degree
$d-1$と内部次数$-a-j$に集中することを用いた。
三角$\mathbf R\tau B\to B\to\mathbf R\omega\pi B$から
\eqref{am:eq:line-hom}を得る。

自然な底変更写像についても確認する。
$M_\ell$の有限階数$B$-自由分解を、syzygyが各内部次数で$R_N$-自由となるように取る。
固定した内部次数でのHom複体は、各項が有限自由$R_N$-加群である。
その剰余はcohomologyが一つの次数にしかない複体であり、
体上のsplit contractionをnilpotent perturbationによって持ち上げると、
有限自由cohomologyの一項複体への実際のhomotopy同値を得る。
したがってこの複体はK-flatであり、通常のtensorが導来tensorを計算する。
項ごとのHomとtensorの同定、filtered colimit、上のsection三角を通じて
自然な底変更写像そのものが準同型であることが従う。
filtered colimitはK-flat複体をK-flat複体へ送るので、ここでunboundedな自由複体の
K-flat性を無条件に用いてはいない。

最後に$\mathbf R\omega\mathcal O(j)$の正のcohomologyは内部次数
$t\le-j-a$にしかない。
$i=1-a>-j-a$なので、$j\ge0$では内部切断後に全て消える。
次数零のsectionは$B(j)$であるから\eqref{am:eq:window-sections}を得る。
\end{proof}

\begin{lemma}[線束生成部分の有限模型]\label{am:lem:linear-category}
\[
 E=\bigoplus_{j=0}^{d-1}\mathcal O(j),\qquad
 \mathcal C_B=\operatorname{thick}_{R_N}\{\mathcal O(j):j\in\ZZ\}
\]
とすると$\mathcal C_B=\operatorname{thick}_{R_N}(E)$である。
全ての対象対のHomは$R_N$-perfectであり、Artin商への底変更と両立する。
例外成分$\langle\mathcal O,\ldots,\mathcal O(a-1)\rangle$を除く右直交を
\[
 \mathcal K_B^{\rm lin}
 =\mathcal C_B\cap\langle\mathcal O,\ldots,\mathcal O(a-1)\rangle^\perp
\]
とすると、有限mutationの射影$p$により$G=p(E)$がこれを生成する。
\end{lemma}

\begin{proof}
式\eqref{am:eq:relative-koszul}を$B$へ導来tensorする。
そのcohomologyは$R_N$と$R_N(-3)$のshiftだけであり、いずれも射影商で零となる。
従って線束の長さ$d$のKoszul型複体は射影商で完全である。
両端の係数の階数が一であるため、任意の連続した$d$個のtwistから全てのtwistを
前後に生成できる。
式\eqref{am:eq:line-hom}と有限cone・直和因子によりproper性と底変更が従う。
同式から$\mathcal O,\ldots,\mathcal O(a-1)$は$R_N$-例外列であり、
Homおよびその$R_N$-双対を用いて左右随伴が作れる。
従って有限個のevaluation coneが射影$p$を与え、$p(E)$が右直交を生成する。
\end{proof}

\subsection{MFのHomのproper性と生成子の持上げ}

\begin{lemma}[nilpotent perturbationによるMFの底変更]\label{am:lem:mf-hom}
有限自由graded MFの任意の二対象$P,Q$に対し、
$\operatorname{Hom}_{\MF}(P,Q)$は$R_N$-perfectである。
自然なHomの導来底変更は通常のMFの底変更Homと準同型である。
\end{lemma}

\begin{proof}
周期変数$t$にcohomological degree $2$、内部次数$-3$を与え、
代数$A$をcohomological degree零に置く。
この二重次数で内部次数零のMF Homを書けば、
各cohomological degreeは有限個の$A_j$の直和であり、有限自由$R_N$-加群である。
負のcohomological方向では必要な内部次数が負となるため零である。
剰余のHom複体$C_0$は、滑らかな古典cubicのgraded MFにおけるHomであり、
通常のOrlov同値により有限次元の有界cohomologyを持つ。

体上で$C_0$から有限cohomology複体$H$へのsplit contractionを選ぶ。
各次数の基底を持ち上げて、元のHom微分を$d_0+\delta$、$\delta\in h\operatorname{End}$と書く。
contractionのhomotopyを$s$とすると、
$(1-\delta s)^{-1}$および$(1-s\delta)^{-1}$は$h^{N+1}=0$により有限和である。
通常のperturbation恒等式をこの有限和に適用すると、
元のHom複体と$H\otimes_kR_N$上の変形微分との間のhomotopy同値を得る。
後者は有界有限自由複体である。
この同値からK-flat性が従い、項ごとのMF Homの底変更が導来底変更と一致する。
\end{proof}

\begin{lemma}[Artin基底上の生成子の持上げ]\label{am:lem:generator}
$P\in\MF^{\ZZ}(A,\Omega)$の剰余$\overline P$が古典graded MF圏を生成するとする。
このとき$P$は$\MF^{\ZZ}(A,\Omega)$を生成する。
\end{lemma}

\begin{proof}
補題\ref{am:lem:mf-hom}により、カテゴリーの底変更から古典MFへの関手は
Hom上で完全忠実である。像は$\overline P$を含むので、冪等完備化後にはMorita同値となる。
生成性を正確に持ち上げるため、$\mathcal D=\operatorname{Ind}(\MF^{\ZZ}(A,\Omega))$へ移る。
$P$はcompactであり、$\operatorname{RHom}(P,-)$はcolimitと可換である。
従って$M\in\mathcal D$が$P$と右直交なら
\[
 \operatorname{RHom}(\overline P,M\otimes_{R_N}^{\mathbf L}k)
 \simeq\operatorname{RHom}(P,M)\otimes_{R_N}^{\mathbf L}k=0.
\]
古典fiberでの生成性から$M\otimes_{R_N}^{\mathbf L}k=0$となる。
$R_N$はその$h$-filtrationによって$D(R_N)$内で$k$の有限extensionから作られる。
従って$R_N\in\operatorname{thick}(k)$であり、$M\otimes^{\mathbf L}k=0$は$M=0$を含意する。
よって$P$は$\mathcal D$のcompact generatorであり、compact部分を取って主張を得る。
ここではsmallなproper圏における右直交の消滅だけから生成性を結論していない。
\end{proof}

\subsection{有限windowによるenhanced比較}

\begin{theorem}[Artin線束残余pairのMF表示]\label{am:thm:comparison}
本節の入力に対し、有限window
\begin{equation}\label{am:eq:window-functor}
 \Psi=\operatorname{stab}\circ
       \bigl((\mathbf R\omega -)_{\ge\,4-d}\bigr)
\end{equation}
は$\mathcal K_B^{\rm lin}$上でdg Morita同値
\begin{equation}\label{am:eq:residual-mf}
 \Psi:\mathcal K_B^{\rm lin}\xrightarrow{\sim}\MF^{\ZZ}(A,\Omega)
\end{equation}
を与える。これはArtin商への底変更と両立する。
さらに$\rho=L_{\mathcal O}\circ(1)$に対し自然に
\begin{equation}\label{am:eq:rotation}
 \Psi\rho\simeq(1)\Psi
\end{equation}
である。特に$\rho$は自己同値となり、MFの三次周期性は
$\rho^3\simeq[2]$を与える。
\end{theorem}

\begin{proof}
まず\eqref{am:eq:window-functor}をdg quasi-functorとして定義する。
derived sectionと内部切断はenhanced関手であり、
stabilizationは補題\ref{am:lem:stabilization}で構成したdg Morita同値の逆である。
この関手が$\mathcal C_B$上で定義されることを生成子で確認する。
補題\ref{am:lem:line-bundles}から、$0\le j\le d-1$について
\[
 (\mathbf R\omega\mathcal O(j))_{\ge i}=B(j)_{\ge i}.
\]
これは$A$上perfectである。
実際、$B(j)$は$A$上長さ一の分解を持ち、有限長商
$T_j=B(j)_{<i}$の次数filtrationの各商は有限自由$R_N$のshiftである。
$R_N$は式\eqref{am:eq:relative-koszul}により$A$上perfectなので、$T_j$も同様である。
よって短完全列
\[
 0\to B(j)_{\ge i}\to B(j)\to T_j\to0
\]
の三項はいずれも$\mathcal D_{B/A}$に入る。
有限coneと直和因子へ延長して$\mathcal C_B$全体に$\Psi$を定義できる。

$0\le j<a$なら$-j\ge i$であり、$B(j)_{\ge i}=B(j)$は$B$上perfectである。
従って$\Psi$は例外成分$\mathcal O,\ldots,\mathcal O(a-1)$を零へ送る。
残る三つの$j=a,a+1,a+2$では
\begin{equation}\label{am:eq:three-images}
 \Psi(\mathcal O(j))\simeq\operatorname{stab}(T_j)[-1],\qquad
 \operatorname{gr}(T_j)=
 \bigoplus_{r=0}^{j-a}B_r\otimes_{R_N}R_N(j-r).
\end{equation}
従って三つの像は、三角的に
$\operatorname{stab}(R_N(a)),\operatorname{stab}(R_N(a+1)),
\operatorname{stab}(R_N(a+2))$を生成する。

剰余上では$\Psi$は通常のOrlovのsection--window関手である。
window $i=1-a$のkernelは$\langle\mathcal O,\ldots,\mathcal O(a-1)\rangle$であり、
その右直交への制限は古典Kuznetsov成分とgraded MFの同値である
\cite[Theorem 6.4 and its proof]{LP-BS}。
右直交と、同定されたwindowの反対側の直交の向きは混同しない。
両者は例外成分による同じVerdier商を表示し、$\Psi$の制限はその商関手である。

全ての自然なHom写像を剰余へ移すと、この古典同値のHom写像となる。
sourceのHomのproper性・底変更は補題\ref{am:lem:linear-category}、
targetは補題\ref{am:lem:mf-hom}による。
$\Psi$の底変更自然性はsectionの計算、内部切断、
補題\ref{am:lem:stabilization}の余核構成から従う。
従ってHom写像のconeは$R_N$-perfectで剰余が零であり、導来Nakayamaにより零である。
これで$\Psi|_{\mathcal K_B^{\rm lin}}$の完全忠実性を得る。
$G=p(E)$の像の剰余は古典MFのgeneratorである。
補題\ref{am:lem:generator}により$\Psi(G)$はArtin上でも全MFを生成するので、
本質的全射性が従う。

最後にrotationの自然性を確認する。
まず$\rho$は$\mathcal K_B^{\rm lin}$を保つ。
実際、$1\le j<a$では
$\operatorname{RHom}(\mathcal O(j),M(1))
=\operatorname{RHom}(\mathcal O(j-1),M)=0$であり、
式\eqref{am:eq:line-hom}より$\operatorname{RHom}(\mathcal O(j),\mathcal O)=0$である。
従ってevaluationの三角から$j\ge1$の直交性が従い、$j=0$の直交性はleft mutationの定義による。
$M\in\mathcal K_B^{\rm lin}$なら
\[
 (\mathbf R\omega M)_i
 \simeq\operatorname{RHom}(\mathcal O(-i),M)
 =\operatorname{RHom}(\mathcal O(a-1),M)=0.
\]
従って内部切断の自然な三角から
\[
 ((\mathbf R\omega M)(1))_{\ge i}
 =(\mathbf R\omega M)_{\ge i+1}(1)
 \simeq(\mathbf R\omega M)_{\ge i}(1)
\]
を得る。stabilizationは内部shiftを保つので、
$\Psi(M(1))\simeq\Psi(M)(1)$となる。
また$\Psi(\mathcal O)=0$であるから、evaluationの三角を$\Psi$で送ると
$\Psi(L_{\mathcal O}M(1))\simeq\Psi(M(1))$を得る。
これらの自然同型を合成して\eqref{am:eq:rotation}を得る。
ここでの周期同型は、指定された中心元のMFにおける周期変数$t$から移したものである。
\end{proof}


\section{正則局所基底からのArtin底変更}\label{ar:sec:regular}

$(R,\mathfrak m,\mathbb C)$を有限生成複素代数の局所化である正則局所環とする。
$A_R=\bigoplus_{j\ge0}(A_R)_j$は中央$R$-代数、$(A_R)_0=R$、
各次数有限自由、剰余が$d\ge5$変数の多項式環であるとする。
$\Omega_R\in Z(A_R)_3$の剰余は滑らかな三次式とし、
$B_R=A_R/(\Omega_R)$、$a=d-3$と置く。
$E_R=\bigoplus_{j=0}^{d-1}\mathcal O(j)$の例外線束成分を除く射影を$G_R$とし、
$\Gamma_R=\operatorname{REnd}(G_R)$と定める。
以下ではこの定義が可能になる有限性と、必要な底変更だけを確認する。

\subsection{局所模型の線束Hom}

\begin{lemma}[正則局所模型に必要な有限性]\label{ar:lem:regular-hom}
$A_R,B_R$は左右noetherianで、$(B_R)_j$は有限自由である。
線束のHomは$R$-perfectであり、
\[
 \operatorname{Ext}^q(\mathcal O(r),\mathcal O(s))=
 \begin{cases}(B_R)_{s-r}&q=0,\\
 (B_R)_{r-s-a}^{\vee}&q=d-2,\\
 0&\text{その他}
 \end{cases}
\]
となる。例外列$\mathcal O,\ldots,\mathcal O(a-1)$による射影が存在し、
$G_R$は線束生成残余部分を生成する。
固定次数のtorsion計算に現れる自由Hom複体は有限自由cohomologyに
homotopy同値であり、任意の通常のtensorで導来tensorを計算できる。
\end{lemma}
\begin{proof}
まずnoetherian性を次数別に確認する。
$A_R$の$\mathfrak m$進随伴次数環は、内部次数ごとの自由性より
$\operatorname{gr}_{\mathfrak m}R\otimes_{\mathbb C}(A_R/\mathfrak m A_R)$であり、
noetherianである。斉次右イデアル$I$の初期形式イデアルの有限生成元を持ち上げ、
それらが生成する右イデアルを$J$とする。
最低$\mathfrak m$次数を順に消すことで、各内部次数$j$と各$q$について
$I_j\subset J_j+\mathfrak m^q(A_R)_j$を得る。
$(A_R)_j/J_j$は有限$R$-加群なので、Krull交叉定理より$I_j=J_j$である。
一般の右イデアルには内部次数filtrationを用い、最高次数の有限降下を行う。
左側も同じであり、$B_R$は商なのでnoetherianである。

第\ref{am:sec:comparison}節の次数別分裂する持上げを、
$R$の$A_R$-加群としてのKoszul分解と、$A_R$の対角Koszul分解に適用する。
前者は項$A_R(-j)^{\binom dj}$、後者は項$(A_R^e)(-j)^{\binom dj}$を持つ。
対角分解を任意の加群へtensorし、$R$の有限大域次元を使えば
$A_R$も有限大域次元を持つ。
また$\Omega_R$による次数別の行列は剰余上単射なので、
単元小行列式により分裂単射であり、各$(B_R)_j$は自由である。

Koszul分解の双対に単元行列成分の消去を適用し、正則中心元のchange-of-ringsを用いると
\[
 \mathbf R\underline{\operatorname{Hom}}_{A_R}(R,A_R)=R(d)[-d],\qquad
 \mathbf R\underline{\operatorname{Hom}}_{B_R}(R,B_R)=R(a)[-(d-1)]
\]
を得る。左右のGorenstein型と$R$-flatnessにより、section--torsion三角に必要な
$\chi$と有限cohomological dimensionを使える
\cite[Section~4, Lemmas~6.1--6.2]{LP-BS}。
$M_\ell=B_R/(B_R)_{\ge\ell}$のfiltrationの商は$(B_R)_j\otimes_R R(-j)$である。
上の双対性から固定内部次数$t$では$j=-t-a$だけが寄与し、
$\ell>j$で安定する。従って
$H^{d-1}(\mathbf R\tau B_R)_t=(B_R)_{-t-a}^{\vee}$、他は零となる。
section--torsion三角から表示した線束Hom公式を得る。
$a>0$なので、固定した内部次数で二つのcohomology群が同時に非零になることはない。

底変更については、$M_\ell$の自由分解を各次数のsyzygyが$R$-自由となるように取る。
固定次数の自由Hom複体は下に有界で、cohomologyは有限自由加群の一つの次数に集中する。
$R$の有限大域次元を用い、完全な射影複体のcycleにdimension shiftingを行うと、
全cycleが射影となり、その複体は分裂する。
従ってcohomology複体へのhomotopy同値があり、通常のtensorが導来tensorを計算する。
項別Hom、filtered colimit、section三角を通じて自然な底変更を得る。

最後にKoszul型の線束完全複体の両端は階数一なので、
連続した$d$本から全twistを前後に生成できる。
Hom公式は例外列とその左右随伴を与え、有限evaluation coneによる射影$p_R$が存在する。
この射影を$E_R$へ適用した$G_R$が線束生成残余部分を生成する。
\end{proof}

\subsection{残余dg代数の底変更}

\begin{proposition}[既存の$\Gamma_R$模型のArtin底変更]\label{am:prop:Gamma-basechange}
この節の正則局所基底$R$上の入力から、
剰余体が$\CC$である局所準同型$R\to R_N$によって$A,B$を底変更したとする。
このとき自然に
\[
 \operatorname{Perf}(\Gamma_R\otimes_R^{\mathbf L}R_N)
 \simeq\mathcal K_{B\otimes_RR_N}^{\rm lin}
 \simeq\MF^{\ZZ}(A\otimes_RR_N,\Omega\otimes1)
\]
であり、mutationで定めたrotation候補はMFの内部shiftへ移る。
\end{proposition}

\begin{proof}
正則な変更先についての底変更だけでは十分でないため、
線束に必要な底変更をArtin基底上で確認する。
正則$R$上で$M_\ell=B/B_{\ge\ell}$の自由分解から作った、固定内部次数のHom複体は、
補題\ref{ar:lem:regular-hom}の証明により有限自由cohomologyへhomotopy同値である。
従って$R_N$への通常のtensorが導来tensorを計算する。
syzygyの次数別$R$-分裂により、底変更後の自由分解も完全である。
そのHom、filtered colimit、section三角を用いると、
変更先について補題\ref{am:lem:line-bundles}で計算したHomと自然に一致する。
よって線束間のHomとその積について必要な底変更を得る。

$R$は正則なので、$-\otimes_R^{\mathbf L}R_N$は有界複体に対して有限Tor振幅を持ち、
射影商への関手が定義される。
線束間のHomの底変更を有限cone・直和因子・evaluationへ延長すれば、
$E_R,G_R$とそのendomorphism dg代数の底変更を得る。
これが第一の同値である。第二の同値とrotationの記述は
定理\ref{am:thm:comparison}による。
この証明は、$D^b(\operatorname{qgr}B)$全体の任意の底変更についての主張を含まない。
\end{proof}


\section{別途与えられたMF同値の残余圏への移送}\label{ar:sec:application}

本節では、非可換MF圏と可換MF圏の比較そのものを構成しない。
その比較を明示的な入力にして、前節までの結果の適用を述べる。

\begin{corollary}[互換なMF比較が与えられた場合]\label{ar:cor:conditional}
$d\ge5$とし、各$N$で第\ref{am:sec:comparison}節の入力$(A_N,\Omega_N)$を取る。
滑らかな特殊ファイバーを持つ形式三次式$W_h\in\mathbb C[[h]][z_1,\ldots,z_d]_3$と、
内部shift・三次周期・Artin商への底変更を保つ同値
\[
 \MF^{\ZZ}(A_N,\Omega_N)\simeq
 \MF^{\ZZ}(R_N[z_1,\ldots,z_d],W_h\bmod h^{N+1})
\]
が各$N$について与えられ、それらが互換であると仮定する。
$B_N=A_N/(\Omega_N)$、$Y_N=V(W_h\bmod h^{N+1})$とすると、互換なpair同値
\[
 (\mathcal K_{B_N}^{\rm lin},\rho)\simeq
 (\Ku_{\rm perf}(Y_N/R_N),\rho_{Y_N})
\]
を得る。ここで右辺は$\operatorname{Perf}(Y_N)$から
$\mathcal O,\ldots,\mathcal O(d-4)$を除いた右直交である。
\end{corollary}
\begin{proof}
定理\ref{am:thm:comparison}を両側に適用し、その間に仮定したMF同値を挟む。
可換側では、線束$\mathcal O(-j)$（$0\le j\le d-1$）の直和が
$D_{\rm QCoh}(Y_N)$をcompactに生成する\cite[Tag~0BQT]{AMStacks}。
これは射影空間の相対対角分解と、閉埋込みのpushforwardの保守性・随伴による。
qcqsスキームではcompact対象とperfect対象が一致する\cite[Tag~09M1]{AMStacks}ので、
全線束の厚い閉包は$\operatorname{Perf}(Y_N)$である。
各比較のrotationと底変更の両立性が、所要の互換性を与える。
\end{proof}

この系は、$W_h$やMF同値の存在証明を代用しない。
非零の複素パラメータへの代数化、完備底での生成性、
全coherent圏との同定についても別の議論が必要である。

\begin{example}[絶対特異圏を代用できない理由]\label{cg:ex:artin-singularity}
$D=\mathbb C[\varepsilon]/(\varepsilon^2)$、$A=D[x]$、$W=x$とする。
有限自由$\mathbb Z/2$-graded MF行列$Q$について$Q^2=xI$を微分すると
$Q(\partial_xQ)+(\partial_xQ)Q=I$となるので、このMF圏は可縮である。
他方、$A/(W)=D$の剰余体は$\varepsilon$による無限周期自由分解を持ち、
その自己Extは任意の非負次数で非零である。
従って剰余体はperfectでなく、絶対特異圏$D_{\rm sg}(D)$は零でない。
周囲でperfectという条件を外した比較は、この例で既に成立しない。
\end{example}

% 2026-11-10 / v22
\section{全てのArtin切断のsmoothnessから完備基底のsmoothnessは従わない}
\label{as:sec:obstruction}

本節では、有限Artin基底での生成子の持上げと、
完備局所基底での生成性の間にある障害を明示する。
出発点となる$\mathbb P^1$上の生成対象の変形例は、
SashaのMathOverflow回答\cite{AMSasha}にある既知の例である。
以下の四次元strict dg模型、および全Artin切断の相対smoothnessと
generic fiberの非smoothnessを併記した整理の先行研究上の初出は未確認である。

この例は、非可換cubicから構成した本稿の$\Gamma_R$が非smoothであることを示すものではない。
示すのは、底上proper、滑らかな特殊ファイバー、
全ての有限Artin切断の相対smoothnessという抽象的条件だけでは、
完備基底上の相対smoothnessを結論できないという事実である。

\subsection{生成対象がgeneric fiberで生成しなくなる族}

$R=\CC[[h]]$、$K=\CC((h))$、$X=\mathbb P^1_R$とし、
$H^0(X,\mathcal O_X(1))$の基底を$s,t$とする。
ここで$s,t$は射影直線の一次sectionであり、
MFの周期変数$t$とは無関係である。
混同を避けるため、本節でのみこの記号を用いる。
二項perfect複体
\begin{equation}\label{as:eq:P}
 P_h=\bigl[\mathcal O_X(-1)\xrightarrow{\,hs\,}\mathcal O_X\bigr]
\end{equation}
をcohomological degrees $-1,0$に置く。

特殊ファイバーでは微分が零なので
\[
 P_0\simeq\mathcal O_{\mathbb P^1}\oplus
                  \mathcal O_{\mathbb P^1}(-1)[1].
\]
これは$\operatorname{Perf}(\mathbb P^1)$の古典的生成対象である。
一方、generic fiberでは$h$は可逆なので、
$s$の零点$p$について
\[
 P_h\otimes_R K\simeq\mathcal O_p.
\]
$\mathcal O_p$から有限coneと直和因子によって作る対象のsupportは$p$に含まれるため、
$\mathcal O_p$は$\operatorname{Perf}(\mathbb P^1_K)$を生成しない。
従って$P_h$も$\operatorname{Perf}(\mathbb P^1_R)$を生成しない。
実際、生成していれば、その有限生成操作を$K$へ底変更することで
generic fiberでも生成してしまう。

\begin{lemma}[各有限切断では生成する]\label{as:lem:artin-generation}
$R_N=R/(h^{N+1})$とし、$P_N=P_h\otimes_RR_N$とする。
このとき$P_N$は$\operatorname{Perf}(\mathbb P^1_{R_N})$を生成する。
\end{lemma}
\begin{proof}
射影直線の相対Beilinson分解により、perfect圏の底変更と
perfect対象間のHomの導来底変更は$R_N\to\CC$と両立する。
従って補題\ref{am:lem:generator}のInd圏における議論を、そのまま
$\operatorname{Perf}(\mathbb P^1_{R_N})$に適用できる。
具体的には、$P_N$と右直交するInd対象$M$について、
$P_N$のcompact性と底変更から
\[
 \operatorname{RHom}(P_0,M\otimes_{R_N}^{\mathbf L}\CC)=0
\]
となる。$P_0$は生成対象なので$M\otimes_{R_N}^{\mathbf L}\CC=0$である。
$R_N$の有限$h$-filtrationにより
$R_N\in\operatorname{thick}_{D(R_N)}(\CC)$だから、$M=0$となる。
これで$P_N$のcompact生成性、従ってperfect圏の古典的生成性を得る。

ここで用いた射影直線の有限模型は、対角の完全列
\[
 0\longrightarrow\mathcal O(-1,-1)
 \longrightarrow\mathcal O
 \longrightarrow\mathcal O_\Delta\longrightarrow0
\]
から得られる。従ってこの底変更の主張は、非正則なArtin基底でも成立する。
\end{proof}

\subsection{四次元strict dg代数}

\begin{proposition}[底上properな明示的dg模型]\label{as:prop:dg-model}
$\Gamma_R=\operatorname{RHom}_X(P_h,P_h)$は、次のstrict unital dg代数で表される：
\begin{equation}\label{as:eq:dg-algebra}
 \Gamma_R^0=Re_-\oplus Re_+,\qquad
 \Gamma_R^1=Rs\oplus Rt,\qquad
 \Gamma_R^q=0\quad(q\ne0,1),
\end{equation}
\[
 d(e_-)=hs,\qquad d(e_+)=-hs,\qquad d(s)=d(t)=0.
\]
$e_-,e_+$は互いに直交する冪等元で、単位は$e_-+e_+$である。
$s,t$は$e_-$から$e_+$へのdegree一の射であり、
\[
 e_+s=s=se_-,\qquad e_+t=t=te_-,
\]
その他の冪等元と射の積、および二本の射同士の積は零である。
特に$\Gamma_R$は$R$上perfectである。
この模型の形成は各$R_N$と$K$への導来底変更と両立する。
\end{proposition}
\begin{proof}
$P_h$の二つの項の間のsheaf Homは
$\mathcal O_X(-1),\mathcal O_X,\mathcal O_X(1)$の直和である。
これらの一次以上のcohomologyは零なので、
大域Homの二項複体が$\operatorname{RHom}$を実際に計算する。
degree $-1$の大域Homは$H^0(\mathcal O_X(-1))=0$であり、
degree零は二つの項のscalar endomorphism、degree一は
$H^0(\mathcal O_X(1))=Rs\oplus Rt$である。
複体の微分を$D=hs$とすると、endomorphismの微分は
$f\mapsto Df-(-1)^{|f|}fD$なので、表示した微分を得る。
積はendomorphismの合成そのものである。

この計算は有限自由$R$-加群とその積だけからなり、
同じline bundle cohomologyの消滅は各底変更後も成立する。
従って通常のtensorがこの模型の導来底変更を計算し、
底変更後の$P_h$のderived endomorphismを与える。
\end{proof}

\begin{lemma}[generic fiberの外積代数模型]\label{as:lem:generic}
$E=K[\epsilon]/(\epsilon^2)$に$|\epsilon|=1$、零微分を与える。
このとき
\[
 E\longrightarrow\Gamma_R\otimes_R K,\qquad
 1\longmapsto e_-+e_+,\qquad \epsilon\longmapsto t
\]
はdg代数のquasi-isomorphismである。
\end{lemma}
\begin{proof}
単位と積と微分を保つことは式\eqref{as:eq:dg-algebra}から直ちに分かる。
$h$が可逆なので
\[
 H^0(\Gamma_R\otimes_R K)=K(e_-+e_+),\qquad
 H^1(\Gamma_R\otimes_R K)=(Ks\oplus Kt)/Ks=K\cdot[t]
\]
であり、他のcohomologyは零である。
従って上の写像は全てのcohomologyで同型となる。
\end{proof}

\subsection{無限derived tensorによる非smoothness}

\begin{lemma}[degree一の外積代数はsmoothでない]
\label{as:lem:exterior-not-smooth}
前補題の$E$は$K$上properだが、homologically smoothではない。
\end{lemma}
\begin{proof}
augmentation $E\to K$によって$K$を左右の$E$-加群とする。
左$E$-加群のsemifree resolution
\[
 Q=\bigoplus_{n\ge0}Eu_n,\qquad |u_n|=0,\qquad
 d(u_0)=0,\quad d(u_n)=\epsilon u_{n-1}\quad(n\ge1)
\]
を考える。$Q\to K$は$u_0\mapsto1$、$u_n\mapsto0$（$n\ge1$）で与える。
$n$に関する増大filtrationはsemifree filtrationであり、
$d^2=0$は$\epsilon^2=0$による。
degree零のcycleは$Ku_0$だけであり、degree一の全ての元
$\epsilon u_n$は$d(u_{n+1})$である。
従って$Q\to K$はquasi-isomorphismである。

右からaugmentation加群$K$をtensorすると
\[
 K\otimes_E^{\mathbf L}K
 \simeq K\otimes_E Q
 =\bigoplus_{n\ge0}Ku_n
\]
となり、右辺の微分は零で、全ての$u_n$はdegree零にある。
従って$H^0(K\otimes_E^{\mathbf L}K)$は無限次元である。
もし$K$が$E$上perfectなら、このderived tensorは
有限個の$K$のshiftから有限coneと直和因子で作られ、
有限次元の有界cohomologyを持つはずなので矛盾する。

さらに、もし$E$の対角が$E^{\rm op}\otimes_KE$上perfectなら、
その有限cell表示の一方の$E$-作用へ$-\otimes_E^{\mathbf L}K$を施すと、
$K$の有限$E$-cell表示を得る。
これは$K$が$K$上perfectであることを用いる。
従ってsmoothnessは$K\in\operatorname{Perf}(E)$を含意するが、
直前でこれを否定した。
よって$E$はhomologically smoothでない。
一方、$E$自体は二次元の有界$K$-複体なのでproperである。
\end{proof}

\begin{theorem}[全Artin切断ではsmoothだが完備基底ではsmoothでない族]
\label{as:thm:obstruction}
命題\ref{as:prop:dg-model}の$\Gamma_R$は、次の全ての性質を持つ。
\begin{enumerate}
 \item $\Gamma_R$は$R$上perfectなstrict dg代数である。
 \item $\Gamma_R\otimes_R^{\mathbf L}\CC$は$\CC$上smoothである。
 \item 各$N\ge0$に対し
 $\Gamma_R\otimes_R^{\mathbf L}R_N$は$R_N$上smoothである。
 \item $\Gamma_R\otimes_R K$は$K$上smoothでなく、
 $\Gamma_R$も$R$上smoothでない。
\end{enumerate}
\end{theorem}
\begin{proof}
第一項は有限自由模型から従う。
補題\ref{as:lem:artin-generation}により$P_N$は
$\operatorname{Perf}(\mathbb P^1_{R_N})$の生成対象であり、
命題\ref{as:prop:dg-model}の底変更から
\[
 \operatorname{Perf}(\Gamma_R\otimes_R^{\mathbf L}R_N)
 \simeq\operatorname{Perf}(\mathbb P^1_{R_N})
\]
となる。
相対射影直線の対角は前述の二項線束分解を持つので、
このperfect圏は$R_N$上smoothである。
smoothnessはdg Morita同値で不変だから、第二・第三項を得る。

補題\ref{as:lem:generic}と
\ref{as:lem:exterior-not-smooth}からgeneric fiberはsmoothでない。
$\Gamma_R$が$R$上smoothなら、そのperfect対角の表示を
flatな底変更$R\to K$へ移してgeneric fiberもsmoothとなるため、矛盾する。
従って第四項も成立する。
\end{proof}

\begin{remark}[この反例の適用範囲]
$R_N\in\operatorname{thick}_{D(R_N)}(\CC)$はArtin基底での生成子の持上げを支える。
しかし$R=\CC[[h]]$は$D(R)$で$\CC$の有限coneと直和因子から作れない。
実際、そのような対象を$K$へtensorすると零になるが、$R\otimes_RK=K\ne0$である。
上の族では、この違いが生成性とsmoothnessの実際の破綻として現れる。

従って、全Artin比較の両立性だけから、完備基底上の生成性や相対smoothnessを補ってはならない。
可換側がsmooth properな射影族であることだけでも、変形した対象の生成性は保証されない。
中心三次非可換商の残余圏へこの原理を適用するには、
当該windowと生成対象についての有限生成証明、
または対角の具体的なperfect表示を別途用いる。
本例は、その特別な主張を否定してはいない。

なお本例の$R$上のcohomologyは
$H^0(\Gamma_R)=R$、$H^1(\Gamma_R)=Rt\oplus(R/hR)s$であり、
cohomologyの底上flatnessは満たさない。
この例が否定する仮定の範囲を、証明していない別の仮定へ広げてはならない。
例\ref{cg:ex:artin-singularity}はMFと絶対特異圏の区別を示すものであり、
本節のsmoothnessの持上げ障害とは別の反例である。
\end{remark}

\begin{thebibliography}{99}
\bibitem{AMChen}
X.-W. Chen, \emph{Module factorizations}, arXiv:2402.11613v3 (27 August 2025).
\url{https://arxiv.org/html/2402.11613v3}.
Theorem~2.10, Corollary~2.11, Theorem~7.2。
余核の三角同値を用い、graded dg比較は本文のexact圏とdg商で構成する。

\bibitem{LP-AZ}
M. Artin and J. J. Zhang, \emph{Noncommutative projective schemes},
Adv. Math. \textbf{109} (1994), 228--287.
\url{https://doi.org/10.1006/aima.1994.1087}.

\bibitem{LP-BS}
J. Burke and G. Stevenson, \emph{The derived category of a graded Gorenstein ring},
arXiv:1507.00830v1 (2015).
\url{https://arxiv.org/abs/1507.00830}.
Section~4, Lemmas~6.1--6.2, Theorem~6.4およびその証明。

\bibitem{AMStacks}
The Stacks Project Authors, \emph{The Stacks Project}.
Tags \href{https://stacks.math.columbia.edu/tag/0BQT}{0BQT}
（有限線束生成子）および
\href{https://stacks.math.columbia.edu/tag/09M1}{09M1}
（compact対象とperfect対象）。

\bibitem{AMSasha}
Sasha, answer to \emph{Is the property of being a dg generator open?},
MathOverflow (5 December 2016).
\url{https://mathoverflow.net/questions/256391/is-the-property-of-being-a-dg-generator-open}.
射影直線上の生成対象がskyscraperへ変形する例を参照する。

\end{thebibliography}
\end{document}
